%==================================================
% SECCIÓN — Structurer.py
%==================================================

\begin{frame}
    \begin{center}
        {\Huge \textbf{Structurer.py}}\\[0.4em]
        {\Large Organización Automática de Texto}
    \end{center}
\end{frame}

%==================================================

\begin{frame}{Objetivo General del Archivo}
El archivo \texttt{structurer.py} transforma texto desordenado
en un resumen organizado automáticamente por temas.

\vspace{0.4em}

\begin{block}{Flujo General}
\begin{center}
Oraciones
$\rightarrow$
Preprocesamiento
$\rightarrow$
Clustering
$\rightarrow$
Scores
$\rightarrow$
Ranking
$\rightarrow$
Resumen Final
\end{center}
\end{block}

\vspace{0.4em}

\textbf{Idea principal:}
\begin{itemize}
    \item Agrupar oraciones similares
    \item Detectar automáticamente el tema
    \item Elegir las frases más importantes
    \item Generar una salida estructurada
\end{itemize}

\end{frame}

%==================================================

\begin{frame}[fragile]{Función \texttt{extract\_topic\_label}}
Esta función descubre automáticamente un nombre representativo
para un grupo de oraciones.

\vspace{0.4em}

\begin{block}{Idea conceptual}
Cluster
$\rightarrow$
Palabras importantes
$\rightarrow$
Mejor palabra
$\rightarrow$
Título
\end{block}

\vspace{0.3em}

\textbf{Código principal:}

\begin{verbatim}
combined = {}

for vec in vectors:
    for word, score in vec.items():
        combined[word] =
        combined.get(word,0) + score
\end{verbatim}

\vspace{0.3em}

Cada oración ``vota'' por ciertas palabras.
La palabra con mayor score domina el tema.

\end{frame}

%==================================================

\begin{frame}{Acumulación de Scores}
El sistema suma la importancia de cada palabra
a través de todas las oraciones del cluster.

\vspace{0.4em}

\begin{block}{Fórmula}
\[
S(w)=\sum_{i=1}^{n} score_i(w)
\]
\end{block}

Donde:

\begin{itemize}
    \item $S(w)$ = score total de la palabra
    \item $score_i(w)$ = score en la oración $i$
    \item $n$ = cantidad de oraciones
\end{itemize}

\vspace{0.3em}

\textbf{Ejemplo:}

\[
\texttt{ia: }2 + 3 = 5
\]

\vspace{0.2em}

La palabra con mayor score será usada como tema principal.

\end{frame}

%==================================================

\begin{frame}[fragile]{Reconstrucción de Palabras}
Los stems no son agradables para mostrar:

\vspace{0.2em}

\begin{center}
\texttt{intelig}
\end{center}

\vspace{0.3em}

Por eso el sistema busca palabras reales similares.

\vspace{0.4em}

\begin{verbatim}
candidates = [
 t for t in all_tokens
 if t.startswith(...)
]
\end{verbatim}

\vspace{0.4em}

\textbf{Ejemplo}

\[
\texttt{intelig}
\rightarrow
\texttt{inteligencia}
\]

\vspace{0.3em}

Luego se selecciona la palabra más larga:

\begin{verbatim}
best = max(candidates, key=len)
\end{verbatim}

\end{frame}

%==================================================

\begin{frame}[fragile]{Construcción de Clusters}
Las oraciones se agrupan según su tema.

\vspace{0.4em}

\begin{verbatim}
clusters.setdefault(label, []).append(i)
\end{verbatim}

\vspace{0.4em}

\textbf{Resultado típico}

\begin{verbatim}
{
 0:[0,1],
 1:[2,3]
}
\end{verbatim}

\vspace{0.4em}

Cada cluster guarda:
\begin{itemize}
    \item oraciones
    \item scores
    \item vectores
\end{itemize}

\vspace{0.3em}

Esto permite construir resúmenes separados por tema.

\end{frame}

%==================================================

\begin{frame}{Importancia del Cluster}
Cada tema obtiene un score promedio.

\vspace{0.4em}

\begin{block}{Promedio}
\[
Promedio =
\frac{\sum score_i}{n}
\]
\end{block}

\vspace{0.4em}

\textbf{Interpretación}

\begin{itemize}
    \item Promedio alto
    $\rightarrow$
    tema importante
    \item Promedio bajo
    $\rightarrow$
    información secundaria
\end{itemize}

\vspace{0.4em}

Luego los clusters se ordenan desde el más relevante
hasta el menos relevante.

\end{frame}

%==================================================

\begin{frame}[fragile]{Ranking de Oraciones}
El sistema selecciona las frases más importantes
de cada tema.

\vspace{0.4em}

\begin{verbatim}
ranked = sorted(
 zip(cluster_scores,
     cluster_sentences),
 reverse=True
)[:top_n]
\end{verbatim}

\vspace{0.4em}

\textbf{Proceso}

\begin{itemize}
    \item Une:
    \[
    (score, sentence)
    \]
    \item Ordena por score
    \item Elige las mejores frases
\end{itemize}

\vspace{0.4em}

Esto funciona como un ranking de buscadores
o sistemas de recomendación.

\end{frame}